\SourceLeaf{287}
\Needspace{6\baselineskip}
\hspace*{\parindent}The logarithms of the sides of the triangle in toises are
\begin{center}
\begin{tabular}{lr}
Hohehagen-Inselsberg&\(4{,}6393865\)\\
Inselsberg-Brocken&\(4{,}7353929\)\\
Brocken-Hohehagen&\(4{,}5502669\)
\end{tabular}
\end{center}
\hspace*{\parindent}The azimuths are
\begin{center}
\begin{tabular}{lrrr}
&\multicolumn{3}{c}{Station Brocken}\\
Inselsberg&\(5^0\)&\(42'\)&\(22''\)\\
Hohehagen&58&49&8\\
&\multicolumn{3}{c}{Station Hohehagen}\\
Brocken&238&9&2\\
Inselsberg&324&23&1\\
&\multicolumn{3}{c}{Station Inselsberg}\\
Hohehagen&144&55&51\\
Brocken&185&35&21
\end{tabular}
\end{center}
\hspace*{\parindent}There is no need here to distinguish between values on the spheroid and those on the sphere. For the logarithms of the distances, a difference first appears in the eighth or ninth decimal place, and for the azimuths, in the thousandths of a second. For our purpose, logarithms to four decimal places and azimuths to the minute are already more precise than necessary. Calculation by the formulas above then gives the following reductions, which are to be added with their signs to the azimuths on the spheroid to obtain the azimuths on the sphere.
\begin{center}
\begin{tabular}{lr}
Brocken-Inselsberg&\(+0''00055\)\\
Brocken-Hohehagen&\(+0{,}00196\)\\
Hohehagen-Brocken&\(-0{,}00238\)\\
Hohehagen-Inselsberg&\(-0{,}00332\)\\
Inselsberg-Hohehagen&\(+0{,}00428\)\\
Inselsberg-Brocken&\(-0{,}00083\)
\end{tabular}
\end{center}
\hspace*{\parindent}Thus the angles of the triangle on the spheroid (between the geodesics) undergo the following changes when reduced to the angles of the spherical triangle (between great-circle arcs)
